\documentclass[11pt,a4paper]{article}
\usepackage[utf8]{inputenc}
\usepackage[T1]{fontenc}
\usepackage[spanish]{babel}
\usepackage{lmodern}
\usepackage{microtype}
\usepackage{enumitem}
\usepackage{tabularx}
\usepackage{booktabs}
\usepackage{hyperref}
\hypersetup{colorlinks=true,linkcolor=blue,urlcolor=blue}
\usepackage{parskip}

\newcommand{\code}[1]{\texttt{#1}}

\title{Auditoria: optimizacion COMPLETA vs PARCIAL}
\author{}
\date{}

\begin{document}
\maketitle

Comparacion a fondo del codigo de la optimizacion completa
(\code{modelo\_completo.py}, \code{optimizar\_completo\_ED.py},
\code{optimizar\_tercer\_mecanismo.py}) contra la optimizacion parcial
(\code{modelo\_simplificado.py}, \code{optimizar\_simplificado\_ED.py},
\code{optimizar\_simplificado\_GA.py}), para verificar que la completa incorpora
las mismas mejoras que la parcial en cuatro ejes:

\begin{itemize}[leftmargin=1.4em]
  \item (a) funcion objetivo Chamfer ponderada por articulacion
  \item (b) penalizacion de monotonicidad (W\_MONO = 5.0)
  \item (c) sintonizacion de hiperparametros con Optuna/TPE (25 trials, semilla 42)
  \item (d) restricciones cinematicas + penalizacion fija de 1000
\end{itemize}

Todas las rutas citadas son relativas a \code{optimizacion/}. Los numeros de
linea corresponden al estado del codigo auditado en la rama
\code{feat/optimizacion-ed-vs-ga}.

\section{Resumen ejecutivo}

La optimizacion completa \textbf{ya esta alineada} con la parcial en los cuatro
ejes. La completa reutiliza exactamente el mismo planteamiento (misma estructura
de funcion objetivo, misma penalizacion de monotonicidad, mismo flujo Optuna/TPE
de dos etapas con las mismas semillas y rangos, y el mismo esquema de
restricciones con penalizacion 1000) y lo \textbf{extiende} de forma coherente
para cubrir la falange distal: agrega el termino de punta (W\_TIP), el perfil
angular DIP (W\_DIP), la monotonicidad de la punta, el tercer mecanismo de cuatro
barras y el tope fisiologico DIP.

\textbf{No se requirio ninguna correccion funcional.} El codigo es consistente
entre ambas versiones y las diferencias que existen son extensiones deliberadas y
documentadas, no inconsistencias. Se verifico por smoke que la
\code{fitness\_function} de ambos modelos devuelve un float finito sobre sus
parametros de referencia (ver el detalle del vector \code{p\_ref} en la seccion
de verificacion smoke, mas abajo).

La tercera etapa que da movimiento a la falange distal se implementa como un
mecanismo de \textbf{cuatro barras}, de forma consistente en todo el codigo
(usuario confirmado).

\section{Comparacion por criterio (resumen)}

\begin{description}[leftmargin=0em,style=nextline]
  \item[(a) Funcion objetivo]\hfill
    \begin{itemize}[leftmargin=1.4em,nosep]
      \item \emph{Parcial (simplificado):} Chamfer ponderada IFP/IFD
            (\code{W\_IFP=0.40}, \code{W\_IFD=0.60}); reg suave
            \code{W\_REG\_DIM=W\_REG\_AUX=5e-4}.
      \item \emph{Completa:} Chamfer ponderada IFP/IFD/punta
            (\code{W\_IFP=0.25}, \code{W\_IFD=0.375}, \code{W\_TIP=0.375});
            ademas perfil DIP \code{W\_DIP=0.3}; misma reg suave.
      \item \emph{Veredicto:} alineado + extendido a la falange distal.
    \end{itemize}
  \item[(b) Monotonicidad]\hfill
    \begin{itemize}[leftmargin=1.4em,nosep]
      \item \emph{Parcial:} \code{monotonicity\_penalty} con \code{W\_MONO=5.0},
            aplicada a IFD.
      \item \emph{Completa:} \code{monotonicity\_penalty} con \code{W\_MONO=5.0},
            aplicada a IFD y punta.
      \item \emph{Veredicto:} alineado + extendido.
    \end{itemize}
  \item[(c) Optuna/TPE]\hfill
    \begin{itemize}[leftmargin=1.4em,nosep]
      \item \emph{Parcial:} 2 etapas, \code{N\_TRIALS\_OPTUNA=25},
            \code{OPTUNA\_SEED=42}, \code{OPT\_SEED=42}; mismos rangos (strategy,
            popsize 15--30, mut\_low, mut\_high, recombination).
      \item \emph{Completa:} idem, mismos trials/semillas/rangos; unica diferencia
            deliberada: \code{polish=False} durante la busqueda corta.
      \item \emph{Veredicto:} alineado (diferencia documentada).
    \end{itemize}
  \item[(d) Restricciones + penalizacion]\hfill
    \begin{itemize}[leftmargin=1.4em,nosep]
      \item \emph{Parcial:} validaciones cinematicas + penalizacion fija 1000 +
            \code{validar\_cinematica()}.
      \item \emph{Completa:} mismas validaciones + penalizacion 1000; agrega
            tercer 4 barras (\code{circle\_intersections}$\to$D3, calibracion
            \code{gamma\_bracket3}) y tope \code{DIP\_MAX\_DEG=35}.
      \item \emph{Veredicto:} alineado + extendido.
    \end{itemize}
\end{description}

\section{Hallazgos detallados}

\subsection{(a) Funcion objetivo}

La estructura de la funcion objetivo es la misma en ambos modelos: Chamfer por
articulacion + penalizacion de monotonicidad + regularizacion suave +
penalizacion fija de 1000 si el mecanismo no ensambla.

Parcial (\code{modelo\_simplificado.py}):

\begin{itemize}[leftmargin=1.4em]
  \item Pesos: \code{W\_IFP = 0.40}, \code{W\_IFD = 0.60}, \code{W\_MONO = 5.0}
        (lineas 72--74).
  \item Regularizacion suave: \code{W\_REG\_DIM = 5e-4}, \code{W\_REG\_AUX = 5e-4}
        (lineas 77--78).
  \item \code{fitness\_function} (lineas 255--277): alinea rigidamente mocap
        contra sim (\code{optimal\_rigid\_transform}), calcula \code{err\_ifp} y
        \code{err\_ifd} con \code{chamfer\_distance}, \code{mono\_ifd} con
        \code{monotonicity\_penalty}, y suma
        \code{shape\_error = W\_IFP*err\_ifp + W\_IFD*err\_ifd},
        \code{mono\_error = W\_MONO*mono\_ifd}, mas \code{reg\_dim} y
        \code{reg\_aux}. Penaliza con \code{1000.0} si \code{run\_kinematics}
        regresa \code{None} (linea 258).
\end{itemize}

Completa (\code{modelo\_completo.py}):

\begin{itemize}[leftmargin=1.4em]
  \item Pesos: \code{W\_IFP = 0.25}, \code{W\_IFD = 0.375}, \code{W\_TIP = 0.375},
        \code{W\_MONO = 5.0}, \code{W\_DIP = 0.3} (lineas 81--85).
  \item Misma regularizacion suave: \code{W\_REG\_DIM = 5e-4},
        \code{W\_REG\_AUX = 5e-4} (lineas 87--88).
  \item \code{fitness\_function} (lineas 250--282): misma alineacion rigida;
        calcula \code{err\_ifp}, \code{err\_ifd} y ademas \code{err\_tip} (punta)
        con \code{chamfer\_distance}; \code{mono\_ifd} y \code{mono\_tip}; un
        termino de error de perfil DIP \code{dip\_error} (lineas 270--271)
        escalado por \code{W\_DIP}; suma
        \code{shape\_error = W\_IFP*err\_ifp + W\_IFD*err\_ifd + W\_TIP*err\_tip},
        \code{mono\_error = W\_MONO*(mono\_ifd + mono\_tip)},
        \code{dip\_pen = W\_DIP*dip\_error}, mas la misma \code{reg\_dim} y
        \code{reg\_aux}. Penaliza con \code{1000.0} si no ensambla (linea 254).
\end{itemize}

\textbf{Veredicto:} consistente. La completa mantiene el planteamiento de la
parcial y lo extiende con el termino de punta y el perfil DIP para cubrir la
falange distal. Los pesos IFP/IFD/punta suman 1.0 en la completa (0.25 + 0.375 +
0.375) y IFP/IFD suman 1.0 en la parcial (0.40 + 0.60), de modo que ambas
reparten la misma unidad de peso de forma coherente.

\subsection{(b) Penalizacion de monotonicidad}

La funcion \code{monotonicity\_penalty} esta definida una sola vez en
\code{comun.py} (lineas 127--137) y la usan ambos modelos con el mismo peso
\code{W\_MONO = 5.0}:

\begin{itemize}[leftmargin=1.4em]
  \item Parcial: \code{mono\_ifd = monotonicity\_penalty(aligned['ifd'])} y
        \code{mono\_error = W\_MONO * mono\_ifd}
        (\code{modelo\_simplificado.py}, lineas 267 y 270).
  \item Completa: \code{mono\_ifd = monotonicity\_penalty(aligned['ifd'])},
        \code{mono\_tip = monotonicity\_penalty(aligned['tip'])} y
        \code{mono\_error = W\_MONO * (mono\_ifd + mono\_tip)}
        (\code{modelo\_completo.py}, lineas 265--266 y 274).
\end{itemize}

Ademas, ambos modelos aplican una restriccion anti-gancho dura sobre el angulo de
la falange medial durante el barrido (\code{delta\_fm < -np.deg2rad(0.5)}
$\to$ \code{None}): \code{modelo\_simplificado.py} lineas 232--233 (dentro de
\code{run\_kinematics}) y \code{modelo\_completo.py} lineas 187--189.

\textbf{Veredicto:} consistente. La completa usa la misma penalizacion y peso, y
la extiende a la punta.

\subsection{(c) Sintonizacion con Optuna/TPE}

El flujo de dos etapas (Optuna/TPE para sintonizar hiperparametros + corrida
profunda) es identico entre \code{optimizar\_simplificado\_ED.py} y
\code{optimizar\_completo\_ED.py}.

Presupuesto y semillas (identicos):

\begin{itemize}[leftmargin=1.4em]
  \item Parcial (\code{optimizar\_simplificado\_ED.py}, lineas 45--49):
        \code{N\_TRIALS\_OPTUNA=25}, \code{DE\_MAXITER\_OPTUNA=40},
        \code{OPT\_MAXITER=300}, \code{OPT\_SEED=42}, \code{OPTUNA\_SEED=42}.
  \item Completa (\code{optimizar\_completo\_ED.py}, lineas 44--48):
        \code{N\_TRIALS\_OPTUNA=25}, \code{DE\_MAXITER\_OPTUNA=60},
        \code{OPT\_MAXITER=600}, \code{OPT\_SEED=42}, \code{OPTUNA\_SEED=42}.
\end{itemize}

Los trials y ambas semillas coinciden (25 trials, semilla 42 para el sampler TPE
y para la ED). El presupuesto por trial y el de la corrida profunda son mayores
en la completa (60 vs 40, 600 vs 300), lo cual es esperable por el mayor tamano
del problema (21 vs 16 parametros); no afecta la simetria del metodo.

Rangos de sugerencia de Optuna (identicos): en ambos \code{objective\_optuna} se
sugieren los mismos hiperparametros con los mismos rangos: \code{strategy} en
\code{['best1bin','rand1bin','best1exp','currenttobest1bin']}, \code{popsize} en
\code{[15, 30]}, \code{mut\_low} en \code{[0.3, 0.7]}, \code{mut\_high} en
\code{[0.8, 1.2]}, \code{recombination} en \code{[0.5, 0.95]}
(\code{optimizar\_simplificado\_ED.py} lineas 73--78;
\code{optimizar\_completo\_ED.py} lineas 78--83).

Sampler y estudio (identicos): ambos usan
\code{optuna.samplers.TPESampler(seed=OPTUNA\_SEED)} y
\code{optuna.create\_study(direction='minimize', ...)}
(\code{optimizar\_simplificado\_ED.py} lineas 90--91;
\code{optimizar\_completo\_ED.py} lineas 95--96).

Diferencia deliberada y documentada: en la completa, la funcion \code{\_correr\_ed}
acepta un parametro \code{polish} (por defecto \code{True}) y
\code{objective\_optuna} lo llama con \code{polish=False} durante la busqueda
corta (\code{optimizar\_completo\_ED.py} lineas 52--66 y 85--86). El docstring
justifica que el pulido local L-BFGS es costoso con 21 parametros y solo hace
falta en la corrida profunda final; la sintonizacion solo necesita comparar
hiperparametros de forma relativa. La corrida profunda si usa \code{polish=True}
(valor por defecto al llamar \code{\_correr\_ed} sin el argumento,
\code{optimizar\_completo\_ED.py} linea 139). En la parcial, la ED usa
\code{polish=True} siempre (\code{optimizar\_simplificado\_ED.py} lineas 53--61).
Esta es una optimizacion de costo, no un cambio de metodo, y esta documentada en
el codigo.

Runner legacy \code{optimizar\_tercer\_mecanismo.py}: este script NO usa Optuna.
Es un runner heredado que fija los hiperparametros a mano:
\code{strategy='best1bin'}, \code{popsize=22}, \code{mutation=(0.5, 1.0)},
\code{recombination=0.7} (\code{optimizar\_tercer\_mecanismo.py} lineas 38--40
para los defaults por entorno y lineas 59--65 para la llamada a
\code{differential\_evolution}). Importa el modelo heredado
\code{exo\_18\_pinza\_fina.py}. La via ``alineada con la parcial'' para la
optimizacion completa es la pareja \code{optimizar\_completo\_ED.py} +
\code{modelo\_completo.py}; \code{optimizar\_tercer\_mecanismo.py} +
\code{exo\_18\_pinza\_fina.py} es el camino legacy que se conserva como
referencia historica.

\textbf{Veredicto:} consistente. La via canonica de la completa
(\code{optimizar\_completo\_ED.py}) usa el mismo Optuna/TPE con los mismos trials,
semillas y rangos que la parcial; la unica diferencia (\code{polish=False} en la
busqueda) es deliberada y esta documentada. El runner legacy sin Optuna existe en
paralelo y no compite con la via canonica.

\subsection{(d) Restricciones cinematicas + penalizacion fija 1000}

Ambos modelos comparten el mismo esquema de restricciones dentro de
\code{run\_kinematics} y \code{validar\_cinematica}:

Validaciones basicas de parametros (regresan \code{None} $\to$ penalizacion 1000):

\begin{itemize}[leftmargin=1.4em]
  \item \code{gear\_ratio <= 0}: \code{modelo\_simplificado.py} linea 103;
        \code{modelo\_completo.py} linea 104.
  \item longitudes minimas \code{min(p[:11]) <= 0.005}:
        \code{modelo\_simplificado.py} linea 105; \code{modelo\_completo.py}
        linea 106.
  \item \code{hsp <= 0 or dsp <= 0}: \code{modelo\_simplificado.py} linea 107;
        \code{modelo\_completo.py} linea 112.
  \item \code{FP\_REAL - 2.0*dsp <= 0.001}: \code{modelo\_simplificado.py}
        linea 109; \code{modelo\_completo.py} linea 114.
\end{itemize}

Restriccion anti-gancho de la medial (\code{delta\_fm >= -0.5 deg}):
\code{modelo\_simplificado.py} lineas 231--233; \code{modelo\_completo.py}
lineas 187--189.

Chequeos de finitud/rango de la trayectoria de salida (\code{abs > 0.3} o no
finito $\to$ \code{None}): \code{modelo\_simplificado.py} lineas 237--239 (IFD);
\code{modelo\_completo.py} lineas 224--226 (punta).

\code{validar\_cinematica()} sobre todo el barrido: comprueba no-ensamble, no
finitud, que se resuelvan las \code{N\_PUNTOS} poses y la monotonicidad de
\code{theta\_fm} (\code{modelo\_simplificado.py} lineas 308--353;
\code{modelo\_completo.py} lineas 318--358).

Penalizacion fija de 1000: \code{return 1000.0} cuando \code{run\_kinematics}
regresa \code{None} en \code{fitness\_function} de ambos
(\code{modelo\_simplificado.py} linea 258; \code{modelo\_completo.py} linea 254).

Extensiones de la completa (falange distal):

\begin{itemize}[leftmargin=1.4em]
  \item Validaciones extra del tercer mecanismo:
        \code{Link9\_3 <= 0.005 or Link10\_3 <= 0.005} y \code{up3\_3 <= 0}
        (\code{modelo\_completo.py} lineas 108--110).
  \item Tercer mecanismo de cuatro barras: calcula el punto de anclaje \code{Pa},
        resuelve \code{D3} por interseccion de circulos
        (\code{circle\_intersections}) y calibra \code{gamma\_bracket3} en la
        primera pose ensamblable (\code{modelo\_completo.py} lineas 208--218). Si
        los circulos no intersectan, regresa \code{None} (linea 211), lo que
        dispara la penalizacion 1000.
  \item Tope fisiologico DIP: \code{DIP\_MAX\_DEG = 35.0} (linea 35) y la
        restriccion \code{if np.ptp(dip\_rel) > np.deg2rad(DIP\_MAX\_DEG): return None}
        (lineas 234--236), que tambien deriva en penalizacion 1000.
\end{itemize}

\textbf{Veredicto:} consistente. La completa hereda todas las restricciones de la
parcial (con la misma penalizacion fija 1000) y agrega las propias del tercer
mecanismo y del tope DIP.

\section{Tercera etapa: mecanismo de cuatro barras (falange distal)}

Al modelo original se le agrego una tercera etapa para dar movimiento a la
falange distal respecto a la proximal, incorporada al analisis cinematico. Esa
tercera etapa se implementa como un mecanismo de \textbf{cuatro barras}, de forma
consistente en todo el codigo:

\begin{itemize}[leftmargin=1.4em]
  \item \code{modelo\_completo.py} linea 9: comentario de encabezado
        \code{ETAPA 3 (4 barras) -> falange distal (articulacion IFD/DIP)}.
  \item \code{modelo\_completo.py} linea 22: \code{Vector de diseno: 21
        parametros (17 del mecanismo base + 4 del tercer 4 barras)}.
  \item \code{modelo\_completo.py} linea 73: comentario \code{\# --- Tercer
        mecanismo de 4 barras (cinematica IFD/DIP) ---}.
  \item La cinematica del tercer mecanismo usa cuatro parametros dimensionales
        (\code{Link9\_3} acoplador, \code{Link10\_3} balancin, \code{back3\_3}
        soporte, \code{up3\_3} standoff dorsal; ver bounds en
        \code{modelo\_completo.py} lineas 74--77 y NOMBRES en lineas 55--56) y
        resuelve el lazo con \code{circle\_intersections}
        (\code{modelo\_completo.py} lineas 209--218).
  \item \code{optimizar\_completo\_ED.py} lineas 114 y 193 tambien lo describen
        como ``tercer 4 barras''.
  \item \code{optimizar\_tercer\_mecanismo.py} lineas 5, 13 y 156 lo describen
        como ``4 barras'' e ``Indices ... tercer mecanismo de 4 barras IFD/DIP''.
  \item \code{exo\_18\_pinza\_fina.py} (modelo legacy) lo describe igual:
        lineas 195, 221, 487, 618 y 670 (``tercer mecanismo de 4 barras''), con
        la misma parametrizacion
        \code{Link9\_3}/\code{Link10\_3}/\code{back3\_3}/\code{up3\_3} y el solver
        \code{\_circle\_intersections} (definido en linea 192, usado en linea 372)
        y la calibracion \code{gamma\_bracket3} (lineas 380--383).
\end{itemize}

\section{Correcciones aplicadas}

\textbf{Ninguna.} No se aplico ningun cambio funcional al codigo.

Justificacion con evidencia:

\begin{itemize}[leftmargin=1.4em]
  \item (a) Funcion objetivo: los pesos y la estructura son consistentes
        (\code{modelo\_simplificado.py} lineas 72--78 y 255--277;
        \code{modelo\_completo.py} lineas 81--88 y 250--282). La completa
        extiende, no contradice.
  \item (b) Monotonicidad: misma funcion y peso \code{W\_MONO=5.0}
        (\code{comun.py} lineas 127--137; \code{modelo\_simplificado.py}
        lineas 267--270; \code{modelo\_completo.py} lineas 265--274).
  \item (c) Optuna: mismos trials (25), mismas semillas (42, 42) y mismos rangos
        de sugerencia (\code{optimizar\_simplificado\_ED.py} lineas 45--49 y
        73--91; \code{optimizar\_completo\_ED.py} lineas 44--48 y 78--96); la
        unica diferencia (\code{polish=False} en la busqueda) es deliberada y
        esta documentada en el docstring de \code{\_correr\_ed}.
  \item (d) Restricciones: mismo conjunto de validaciones y misma penalizacion
        1000 (referencias de linea en la seccion (d)).
\end{itemize}

Verificacion smoke (sin excepcion, float finito) sobre los parametros de
referencia, ejecutada tras la auditoria.

Modelo completo: se usa el vector \code{p\_ref} del bloque \code{main} de
\code{optimizar\_completo\_ED.py} (lineas 117--120), es decir

\begin{verbatim}
p_ref = [0.018, 0.020, 0.035, 0.049, 0.025, 0.020, 0.025,
         0.055, 0.035, 0.052, 0.04601, 0.017, 0.018,
         np.deg2rad(51.39), np.deg2rad(38.78), 2.0, np.deg2rad(109),
         0.025, 0.035, -0.012, 0.002]
\end{verbatim}

Con ese vector, \code{modelo\_completo.fitness\_function(p\_ref)} = 0.036594
(float finito). El valor es reproducible con solo copiar ese \code{p\_ref}.

Modelo simplificado: \code{modelo\_simplificado.fitness\_function(p)} devuelve un
float finito (sin excepcion) para un vector de 16 parametros dentro de los bounds
declarados en \code{modelo\_simplificado.py}; el valor numerico exacto depende
del vector elegido y no se fija aqui. Lo verificable y reproducible es que la
funcion devuelve un float finito. Un vector de prueba valido es el que se usa en
el smoke de la entrega (FEAT-001):

\begin{verbatim}
[0.02, 0.04, 0.03, 0.05, 0.025, 0.02, 0.03, 0.05,
 0.035, 0.05, 0.033, 0.017, 0.008, 0.7, 3.0, -1.4]
\end{verbatim}

Como no se toco ningun \code{.py}, la \code{fitness\_function} de ambos modelos
sigue devolviendo un float finito.

\section{Recomendaciones}

\begin{enumerate}[leftmargin=1.6em]
  \item Mantener \code{optimizar\_completo\_ED.py} como la via canonica de la
        optimizacion completa (es la alineada con la parcial via Optuna).
        Documentar en la entrega que \code{optimizar\_tercer\_mecanismo.py} +
        \code{exo\_18\_pinza\_fina.py} es el camino legacy con hiperparametros
        fijos, conservado como referencia.
  \item Dejar constancia en la redaccion de la tesis/articulo de la diferencia
        deliberada \code{polish=False} durante la sintonizacion de la completa,
        para que la comparacion metodologica quede documentada.
  \item Si en algun momento se desea simetria total del pulido, se podria activar
        \code{polish=True} tambien en la busqueda de la completa, pero no es
        necesario para la validez de la comparacion y encarece la sintonizacion.
\end{enumerate}

\end{document}
